\section{Discussion}

\label{sec:discussion}

% PAPER_STUDIO_PARAGRAPH:D-P1

Table~\ref{tab:main-results} is designed to compare permutation-invariance rates across option-order counts and evaluation criteria, while Figure~\ref{fig:planned-results} is designed to display the distributional trend of these rates across the same dimension. Together, these artifacts test whether sensitivity to option ordering varies systematically with the number of alternatives presented. Because no measurements are yet available, the current evidence gap is explicit: we cannot yet determine whether the hypothesized trend exists, nor whether its direction, if any, is monotonic. We therefore treat the possibility of a systematic ordering effect as a hypothesis only, to be confirmed or rejected once the planned experiments are executed and the placeholders are populated.

% PAPER_STUDIO_PARAGRAPH:D-P2

A primary limitation is that the planned study measures permutation invariance only through aggregate rates across option-order counts, which may obscure item-level variation in sensitivity. The design does not yet isolate whether observed instability stems from genuine order effects or from model variance across repeated samplings, since the evaluation protocol has not fixed a decoding temperature or a number of stochastic passes. A further threat to validity is that the candidate pool generated by reordering options is not equivalent to the set of examples admitted to a refit, so conclusions about selection behavior must remain scoped to the former. Finally, without held-out answer distributions we cannot distinguish a model that flips between plausible alternatives from one that degrades to near-random choice.
